\section{Procedimiento Experimental}

El procedimiento experimental se realizó siguiendo una secuencia ordenada de
mediciones para el resorte metálico y la liga (tira de jebe), considerando
tanto el proceso de carga como de descarga.

\begin{enumerate}

\item \textbf{Mediciones iniciales:}
    Previo al montaje, se determinó la masa de cada material con la balanza,
    su longitud natural con la regla métrica y las dimensiones de su sección
    transversal con el vernier. Los valores obtenidos se reportan en la
    Sección~\ref{sec:resultados}.

\item \textbf{Montaje experimental:}
    Se suspendió el material de estudio verticalmente en el soporte
    universal, asegurando que quedara completamente alineado y en reposo
    antes de realizar cualquier medición. Este montaje se empleó tanto
    para el resorte metálico como para la liga elástica.

\item \textbf{Proceso de carga:}
    Se colocó una primera masa en el extremo libre del resorte y se midieron
    la nueva longitud~$\ell$ y las dimensiones de la sección transversal
    deformada con el vernier. Luego, sin retirar la masa inicial, se
    añadieron sucesivamente las demás masas de forma acumulativa, registrando
    en cada caso la longitud y la sección transversal correspondientes.
    Estos datos permiten calcular tanto el esfuerzo técnico~$\sigma_0$
    como el esfuerzo real~$\sigma = F/S$.

\item \textbf{Proceso de descarga:}
    Una vez colocadas todas las masas, se procedió a retirarlas en orden
    inverso, registrando la longitud y la sección transversal del material
    en cada etapa. Este proceso permitió evaluar posibles efectos de
    histéresis.

\item \textbf{Repetición con la liga:}
    Se repitió el mismo procedimiento experimental utilizando la liga (tira
    de jebe) como cuerpo de estudio. Se tomaron las mismas mediciones de
    longitud y sección transversal en los procesos de carga y descarga,
    cuidando de no aplicar fuerzas que excedan su límite elástico.

\item \textbf{Consideraciones experimentales:}
    Durante todas las mediciones se verificó que el sistema estuviera en
    equilibrio antes de registrar los datos, minimizando errores por
    oscilaciones. Asimismo, se evitó el error de paralaje en la lectura de
    la regla y se comprobó el correcto funcionamiento del vernier antes de
    cada medición.

\end{enumerate}